\begin{chapterbox}
    \chapter{Der Tyrann von Andor (WIP)}
    \label{wip-der-tyrann-von-andor}
    \az{74 bis 75}

    Der tyrannische Regent Ken Dorr hat die Macht über Andor ergriffen. Jarid und Trieest erleiden Verluste und erschöpfen alle Möglichkeiten, zukünftiges Leiden zu verhindern. Sie brauchen dringend Verbündete. Doch Ken hat ein Gespür dafür, Motivationen gegeneinander auszuspielen.

    Inhaltswarnung: Diese Geschichte ist düsterer und tödlicher als meine üblichen. 
\end{chapterbox}



Diese Fan-Geschichte ist noch in Arbeit. Ich teile sie bereits hier, damit sie für mich nicht verloren geht, falls mein Laptop spontan verschwinden sollte, und damit sie für euch nicht völlig verloren geht, falls ich spontan verschwinden sollte. Fühlt euch in letzterem Fall vollkommen frei, hierrin zu schmökern und nach Lust und Laune weiter daran zu werkeln. Aber solange ich ein aktiver Andori bleibe, lest bitte nicht zu genau nach, was ich geplant haben könnte. LG BBB :D



\newpage
\section{Prolog}
\az{Jahr 39}

Ken und seine \gend{Diebesgefährt*innen}{Diebesgefährten} flohen durch das kreaturenverseuchte Rietland, verstoßen von der Rietburg für den Diebstahl des Feuerschilds. 

Eine Skral-Sippe nahm sie gefangen.

Ken konnte sich zu einem Zweikampf gegen einen Skral durchreden.

Er wurde gezwungen, seine Kleidung liegen zu lassen, und ihm waffenlos gegenüberzutreten.

Ken hatte ein klein wenig Gift dabei, das er nutzen konnte, um sein Gegenüber zu besiegen, aber nicht, ehe jenes ihm nicht den Rücken mit scharfen Klauen zerfleischte.

Ken humpelte mehr davon, als er rannte.

Kens \gend{Gefährt*innen}{Gefährten} schrien hinter ihm, als die Skrale sie töteten.

Er schaffte es zur Küste des Wachsamen Waldes.

Dort hielt der Bewahrer Tion Wache. Ken überzeugte ihn davon, ihm eine Überfahrt in den Norden zu verschaffen.

Ken reiste davon, nichts als Hass in seinem Blick und in seinem Sinn. Auf die Kreaturen. Auf Brandur. Auf seine nichtsnutzigen \gend{Diebesgefährt*innen}{Diebesgefährten}. Auf alle, die ihm geschadet hatten. Sie würden allesamt brennen.




\newpage
\section{Kapitel}
\az{Jahr 74}

Thorald jagte den Königswolf Lonas bis zu den Ausläufern des Grauen Gebirges.

Fell und andere Trophäen, die Thoralds Gefolge trug, zeugten von seinem Mord an Merla und Rutan.

Elga erschien auf einer Klippe, bezeugte ihr Überleben und flehte Thorald an, von seiner falschen Jagd abzulassen. Sie verschwand von der Klippe, um einen Weg hintendurch zu Thorald zu finden.

Thorald entspannte sich und begann, gleichzeitig zu lachen, weil es Elga gut ging, und zu weinen, weil Ken ihn wieder einmal manipuliert hatte und er das Königsrudel dezimiert hatte.

Da begann Lonas zu knurren, und raste davon. Er hatte \gend{Gegner*innen}{Gegner} hinter dem Hügel wahrgenommen.

Da, ein Schrei! Elga?

Thorald schreckte auf, hüpfte auf sein Pferd und war darauf und daran, Lonas zu folgen. 

Doch sein Gefolge, allesamt von Ken gekauft, schritten ein, nun, wo Thorald Bescheid wusste und der Königswolf nicht mehr direkt daneben stand, und erdolchten Thorald.

Sie schnappten sich die Rietgraskrone und ritten davon.

Orfen erschien, um Thorald zu warnen, dass die Wölfe unschuldig waren, dass er Beweise für die Rückkehr von menschenfressenden Trollsippen im Rietland hatte. Doch er war zu spät. Thorald kippte tot aus dem Sattel.

Nun vernahm auch Orfen das, was Lonas vernommen hatte. Er schlich sich zum nächsten Hügel, und erblickte die Skelettarmee der Krahder. Er sah auch Käfige mit lebendigen Gefangenen, darunter Elga.

Eine Horde Skelett\gend{krieger*innen}{krieger} entdeckte Orfen, jagte ihm nach und setzte ihn fest. Orfen wehrte sich tapfer. Dies war seine Bergwelt, in der er sich am besten auskannte, in der er am stärksten war. Doch gegen die Übermacht konnte er nicht ankommen. Schließlich trat ein Krahder vor ihn und zerschlug sein Schwert. Es blieb im Matsch zurück, während Orfen davongeschleppt wurde.

Der Krahder sah über die Klippe herunter zu Thoralds Leichnam. Er winkte zwei Skeletten zu, sie sollten auch diesen Leichnam zu den anderen schaffen. Er würde ihrer Sache als eines von Hunderten Skeletten in ihrem immer größer werdenden Heer dienen.

\trennzeichen

Kens Handlanger brachten die Rietgraskrone zurück zu Ken, welcher auf der Rietburg bereits Vorkehrungen für den Fall vorbereitet hatte, dass eine rasche Übernahme der Regierung nötig sein würde.

Ein Handlanger verlangte viel Gold, um sein Schweigen zu bewahren.

Ken zückte eine Arcuballiste aus Werftheim und tötete sie alle.

Dann setzte er die Krone auf und betrachtete sich selbst im Spiegel.




\newpage
\section{Kapitel}

Trieests Sippe wurde von hungrigen \gend{Bäuer*innen}{Bauern} angegriffen.

Calrai wurde getötet, obwohl ein königlicher Schutzbrief die Sippe vor Kopfgeld\gend{jäger*innen}{jägern} hätte schützen sollen.

\trennzeichen

Skral-Häuptling Narb ordnete einen Angriff auf die Rietburg an. 

Trieests Sippe stellte sich Narb im Kampf. Narb nahm die Sache locker, selbst als seine Armee schwächelte.

Er erwartete einen Verrat im Inneren der Burg und dass sich jeden Moment die Tore öffnen würden.

Doch die Tore der Rietburg blieben geschlossen. Und auf einmal dezimierten schwere Ballisten seine Horden – und Trieests.

Narb rief zum geordneten Rückzug auf.

\trennzeichen

Als Narb bei seiner Höhle im Skralberg ankam, erwartete ihn ein schreckliches Bild. 

Viele Frauen und Kinder litten unter Magenbeschwerden oder lagen regungslos am Boden. Überall waren Pfützen des Gebräus. Es stank nach Erbrochenem und nach Schwefel. Die Drachenfässer Rachenputzer waren vergiftet gewesen!

Viele überlebende \gend{Skral*innen}{Skrale} waren unverletzt, aber zutiefst verstört.

Am Ende der Höhle erkannte Narb seinen pfeifenrauchenden Informanten. Die Tore der Rietburg hatten sich nicht geöffnet. Der Freudentrank war vergiftet gewesen. Wenn jemand ihn verraten hatte, dann höchstwahrscheinlich er.

Der Informant humpelte an Narb vorbei, zum Höhleneingang. Er würde ihm nicht entkommen. Narb folgte ihm und riss ihm die Kapuze herunter.

Da richtete sich der Informant auf. Sein scheinbar verletztes Bein trug ihn problemlos. Etwas glänzte auf der Stirn des Informanten. Er trug eine goldene Krone, die Rietgraskrone Andors!

Langsam drehte der Informant sich um und zeigte sein bösartigstes Grinsen. Narb erkannte ihn. Ken Dorr, den Regenten Andors.

Ken bließ Narb Rauch ins Gesicht und schleuderte locker seine Pfeife an ihm vorbei. 

Der Rachenputzer entzündete sich sofort. Feuer loderte wie entlang einer Lunte. Noch nicht angebrochene Fässer explodierten. Die vom Gift verschont gebliebenen \gend{Skral*innen}{Skrale} gingen in Flammen auf.

Zum ersten Mal in seinem Leben war Narb erstarrt. Er konnte sich nicht zu einer Handlung durchringen. 

Ken starrte ihm grinsend ins Gesicht, sich an Narbs Leid ergötzend. Er prägte sich den Augenblick so gut er konnte ein.

Dann schnitt er ihm die Kehle durch.


\newpage
\section{Kapitel}

Trieests Sippe enthüllte Kens Machenschaften gegenüber den \gend{Bewohner*innen}{Bewohnern} der Rietburg.




\newpage
\section{Kapitel}

Jarid und Trieest feierten in der Taverne. Dies war ein wichtiger Anfang für den Kampf gegen Ken.

Trieest verließ die Taverne, um zu pinkeln.

Er kehrte nie zurück.






